%% kapa: exam=1 round=37 year=2026 period=1 form=A subject=civil_law subject_ko=민법 seq=34 number=34
%% source: Q-Net 감정평가사 기출문제 (cst00309) artlSeq=5258924 file=raw/1008/5258924/2026년 제37회 감정평가사 1차 1교시.pdf
%% conversion: kordoc 4.24.51 -> tools/gen1.py (auto); review=proofread-sonnet; flags=none
\begin{question}{34}
민사유치권에 관한 설명으로 옳은 것은? (다툼이 있으면 판례에 따름)
\begin{choices}
  \item 동산유치권은 선의취득의 대상이 될 수 없다.
  \item 피담보채권이 변제기에 이르지 않았다는 사정은 유치권의 행사를 방해하지 않는다.
  \item 유치권에는 우선변제적 효력이 인정되지 않으므로 유치권자는 채권의 변제를 받기 위하여 유치물을 경매할 수 없다.
  \item 저당권의 실행으로 부동산에 경매개시결정의 기입등기가 이루어지기 전에 유치권을 취득한 사람은 경매절차의 매수인에게 유치권으로 대항할 수 없다.
  \item 채무자의 소유에 속하지 않는 물건에 대해서는 유치권이 성립할 수 없다.
\end{choices}
\end{question}
\answer{1}
